\section{Research Directions}
\label{sec:directions}

We group the 26 contributions into five research directions. Several papers touch more than one direction; we place each where its main contribution lies.

\subsection{Grounding Agents in the Repository}

Six papers examine what an agent must know, and where it should look, before it writes a patch. Al Awad and Ivanov~\cite{agdev26:alawad} treat repository exploration as a separately measurable and budgetable stage: cheaper explorer models retain most of the localisation quality of the strongest one while using far fewer tokens, and the right choice depends on how the next stage consumes the candidate files. Mahesh~\cite{agdev26:mahesh} isolates routing before any model call and shows that how issue and path tokens are aligned largely determines ranking quality, giving an auditable, zero-inference baseline. Kerbel et al.~\cite{agdev26:kerbel} version agent-facing project knowledge alongside the code; in a pre-registered study on SWE-bench Verified, context written by a stronger model improved localisation by a small local agent, whereas context it wrote itself did not. Dente et al.~\cite{agdev26:dente} decompose repair into localisation, reproduction and editing under harness control, which lifts small open-weight models substantially and lets each stage be assigned to a different model. Nyberg and Samsten~\cite{agdev26:nyberg} show that an agent's next tool call is predictable by a small schema-conditioned model, even when the tool interface drifts, and Borysenko~\cite{agdev26:borysenko} finds that coding agents retrieving documentation leave stable, distinguishable HTTP signatures.

\subsection{Reliability of Agentic Loops}

Five papers question whether iterating is enough. Gao et al.~\cite{agdev26:gao} show that generate--test--revise loops can lose correct patches, particularly when revisions are guided by stale execution traces, and propose a typed loop contract that binds verifier evidence to exact code states and keeps verified checkpoints. Spiridonov~\cite{agdev26:spiridonov} argues that a workflow's value lies in how well it converts a verifier, such as a test suite, into quality, and finds that verifier use is reliable only when it is a deterministic step rather than the model's choice. Khanna~\cite{agdev26:khanna} replays public trajectories and finds that, under a strict false-abort budget, four inexpensive stop signals stay silent on most doomed runs; the paper proposes a reporting protocol for future detectors. Osmanoglu~\cite{agdev26:osmanoglu} treats planning reliability as a systems problem, combining evidence collection, deterministic validation, fallbacks and human checkpoints, and Sun~\cite{agdev26:sun} offers a protocol for evaluating whether procedural skills learnt from a few debugging trajectories transfer, including negative transfer.

\subsection{Evaluation, Benchmarks and Data}

Three papers target what and how we measure. BDD-Bench~\cite{agdev26:jureczka} evaluates agents on executable behavioural specifications across the chronological history of five repositories; the best of thirteen models resolves 45\% of instances, and agents often pass visible tests yet fail hidden variants of the same behaviour. Zergetaev~\cite{agdev26:zergetaev} shows that popular efficiency plugins cut output tokens, which make up a small share of a multi-turn session's cost, and in correctness-gated runs leave the bill unchanged. Ouajdi et al.~\cite{agdev26:ouajdi} generate tool-use training traces execution-first, synthesising tasks only for workflows that actually run. Methodological care was visible throughout the programme, with pre-registration~\cite{agdev26:kerbel,agdev26:khanna}, repository-clustered uncertainty~\cite{agdev26:alawad,agdev26:mahesh} and run-to-run consistency~\cite{agdev26:le,agdev26:stengg} all reported.

\subsection{Engineering Multi-Agent Systems}

Four papers treat agent systems themselves as software to be modelled and engineered. MOSAICO~\cite{agdev26:vu} provides an infrastructure in which a collaboration agent orchestrates solution, supervision and consensus agents according to a workflow, demonstrated on metamodelling, UML generation and code documentation. Ait et al.~\cite{agdev26:ait} define agentic swarms on the BESSER low-code platform through four cross-referenced models, from which the swarm is generated and deployed. Garcia-Dominguez et al.~\cite{agdev26:garciadominguez} generate domain-specific agents that ``think in code'', allowing local open-weight models to approach a frontier model on UML modelling at a higher token cost. Pesare et al.~\cite{agdev26:pesare} argue that architecture shapes values, proposing patterns for privacy, pluralism and fairness in multi-agent systems.

\subsection{Agents across the Lifecycle}

Eight papers apply agents beyond code generation. For \emph{specification}, Soto et al.~\cite{agdev26:soto} represent specifications as a typed, queryable graph that agents consume at runtime, evaluated in industry with 19 practitioners. For \emph{design}, Kakran et al.~\cite{agdev26:kakran} automate ATAM and SAAM architecture evaluations as drafting aids, and Extender~\cite{agdev26:chen} detects SOLID violations by actually extending the code and inspecting the resulting changes. For \emph{testing}, Teach-to-Crash~\cite{agdev26:ghazal} pairs a teacher and a student LLM to generate collision-inducing driving scenarios, and Sari et al.~\cite{agdev26:sari} annotate test types in machine-learning commits through multi-LLM voting. For \emph{security}, Arasan et al.~\cite{agdev26:arasan} combine agentic harness generation with symbolic execution to find, or rule out, inputs that trigger several bugs at once, and Le et al.~\cite{agdev26:le} show that local LLMs can prioritise vulnerability review of robotics software but should not be allowed to dismiss findings. For \emph{maintenance}, katIA~\cite{agdev26:stengg} refines static change-impact analysis with LLM agents. Across all eight, agents complement rather than replace deterministic analysis and human judgement.
